\SourceLeaf{257}%
und ebenso aus der Formel \(r\cos\psi=\tfrac{1}{2}\mathord{.}\frac{\mathrm{d}rr}{\mathrm{d}q}\)
\[
\begin{aligned}
\relax[3]\quad r\cos\psi={}&q+\tfrac{2}{3}f^0ppq+\tfrac{1}{2}f'p^3q+(\tfrac{2}{5}f''-\tfrac{4}{45}f^0f^0)p^4q\quad\text{etc.}\\
&+\tfrac{3}{4}g^0ppqq+\tfrac{3}{5}g'p^3qq\\
&+(\tfrac{4}{5}h^0-\tfrac{14}{45}f^0f^0)ppq^3
\end{aligned}
\]
Damit ist zugleich der Winkel \(\psi\) bestimmt. Entsprechend entwickelt man zur Berechnung des Winkels \(\varphi\) zweckmäßiger die Reihen für \(r\cos\varphi\) und \(r\sin\varphi\). Dazu dienen die partiellen Differentialgleichungen
\[
\begin{aligned}
\frac{\mathrm{d}\mathord{.}r\cos\varphi}{\mathrm{d}p}&=n\cos\varphi\mathord{.}\sin\psi-r\sin\varphi\mathord{.}\frac{\mathrm{d}\varphi}{\mathrm{d}p}\\
\frac{\mathrm{d}\mathord{.}r\cos\varphi}{\mathrm{d}q}&=\cos\varphi\mathord{.}\cos\psi-r\sin\varphi\mathord{.}\frac{\mathrm{d}\varphi}{\mathrm{d}q}\\
\frac{\mathrm{d}\mathord{.}r\sin\varphi}{\mathrm{d}p}&=n\sin\varphi\mathord{.}\sin\psi+r\cos\varphi\mathord{.}\frac{\mathrm{d}\varphi}{\mathrm{d}p}\\
\frac{\mathrm{d}\mathord{.}r\sin\varphi}{\mathrm{d}q}&=\sin\varphi\mathord{.}\cos\psi+r\cos\varphi\mathord{.}\frac{\mathrm{d}\varphi}{\mathrm{d}q}\\
n\cos\psi\mathord{.}\frac{\mathrm{d}\varphi}{\mathrm{d}q}+\sin\psi\mathord{.}\frac{\mathrm{d}\varphi}{\mathrm{d}p}&=0
\end{aligned}
\]
deren Verbindung die folgenden Beziehungen liefert
\[
\begin{aligned}
\frac{r\sin\psi}{n}\mathord{.}\frac{\mathrm{d}\mathord{.}r\cos\varphi}{\mathrm{d}p}+r\cos\psi\mathord{.}\frac{\mathrm{d}\mathord{.}r\cos\varphi}{\mathrm{d}q}&=r\cos\varphi\\
\frac{r\sin\psi}{n}\mathord{.}\frac{\mathrm{d}\mathord{.}r\sin\varphi}{\mathrm{d}p}+r\cos\psi\mathord{.}\frac{\mathrm{d}\mathord{.}r\sin\varphi}{\mathrm{d}q}&=r\sin\varphi
\end{aligned}
\]
Hieraus lassen sich die Reihen für \(r\cos\varphi,r\sin\varphi\) leicht entwickeln. Ihre ersten Glieder müssen offenbar \(p\) beziehungsweise \(q\) sein. Man erhält
\[
\begin{aligned}
\relax[4]\quad r\cos\varphi={}&p+\tfrac{2}{3}f^0pqq+\tfrac{5}{12}f'ppqq+(\tfrac{3}{10}f''-\tfrac{8}{45}f^0f^0)p^3qq\quad\text{etc.}\\
&+\tfrac{1}{2}g^0pq^3+\tfrac{7}{20}g'ppq^3\\
&+(\tfrac{2}{5}h^0-\tfrac{7}{45}f^0f^0)pq^4
\end{aligned}
\]
\[
\begin{aligned}
\relax[5]\quad r\sin\varphi={}&q-\tfrac{1}{3}f^0ppq-\tfrac{1}{6}f'p^3q-(\tfrac{1}{10}f''-\tfrac{7}{90}f^0f^0)p^4q\quad\text{etc.}\\
&-\tfrac{1}{4}g^0ppqq-\tfrac{3}{20}g'p^3qq\\
&-(\tfrac{1}{5}h^0+\tfrac{13}{90}f^0f^0)ppq^3
\end{aligned}
\]
Durch Verbindung der Gleichungen [2], [3], [4], [5] ließe sich eine Reihe für \(rr\cos(\psi+\varphi)\) herleiten, daraus durch Division durch die Reihe [1] eine Reihe für \(\cos(\psi+\varphi)\) und aus dieser eine Reihe für den Winkel \(\psi+\varphi\) selbst. Eleganter erhält man dieselbe Reihe jedoch
